\section{Vox Populi}

\begin{quotex}
And those people should not be listened to who keep saying the voice of the people is the voice of God, since the riotousness of the crowd is always very close to madness.
\flright{\textsc{Alcuin} in a letter to Charlemagne}
\end{quotex}


As the worldwide beacon of light for the virtues of democracy, the recent elections in the USA deserve an analysis by the standards of the historiography of the Right. At this point, it must be regarded as an experiment, subject to future development and modification. The lineaments of a Traditional organization must be more fully developed in order to gauge the full extent of the deviation.

The first requirement is to stay firm to idealism and objectivity. That is, the ideals of the Brahman and Kshatriya are the standards, which must be applied objectively, free from partisanship, passion, or attachment. This is not the same as an alleged ``scientific'' objectivity that relies heavily on statistics, reducing everything to quantity. In order to reach into the roots of thought, i.e., the qualitative aspect, we will need sometimes to rely on what is derided as ``anecdotal'' evidence.

Although native-born Americans have no interest beyond the borders, the rest of the world is always interested in the American elections. Since we now have a large stock of aliens, I have been asking about their preference. The most interesting was from a Chinese fellow who preferred Obama because he felt Obama could relate more easily to non-Americans and their culture (this was the view, according to him, of his countrymen). Since the rest of the world, apart from Israel, is overwhelmingly pro-Obama, this point of view must be quite prevalent. The corollary of how they regard the white majority in the USA can be deduced by the reader.

To be objective means also to live without illusion and to disregard all pretense; note, however, that this differs from mere cynicism. The first task is to determine the goal and method of an election. Ostensibly, the purpose is to select the supreme leader of the temporal power---called the President---of the country. However, only God can select the legitimate leader, who is anointed by the spiritual authority, and the \emph{vox populi} can never be the \emph{vox Dei}. At best, the people can select a ``manager'' or ``chief executive'' who can run the affairs of state for a set period, but without any supernatural sanction.

The Republicans, as the party of the bourgeoisie, seem to accept this, given their choice of Romney who would have run the country like a business. Nevertheless, the Republicans wistfully wished instead for a charismatic candidate as the Democrats had. The latter, as the party of the serfs, appealed to material interests in a very emotional way. The major difference, and this is almost too obvious to mention, is that the one party appeals to those who want to hold onto their money, and the other party to those who want to confiscate the money of the former. This has nothing to do with justice.

The method of an election is to count votes, on the basis of ``one man, one vote''. As Rene Guenon points out, man reduces himself to the ``condition of numerical units''. Each candidate promises unity, but to quote Guenon again, there are only men ``parodying unity, yet lost in the uniformity and indistinction of the mass, that is, in pure multiplicity and nothing else''. In this system, the sage and the idiot are indistinguishable, yet this method is considered absolutely just.

In this indistinction, any talk about ``issues'' or the ``best man for the job'' are mere pretenses, as quality cannot be reduced to quantity. The populace divides along predictable lines, picking sides with the same amount of wild enthusiasm and lack of real thought as the Roman people choosing the blue or the green chariot racing team. People even pray for their candidate to win, but since the election is based solely on quantity, that makes as much sense as praying for 2+2 to no longer be 4.

Since the voice of the people is not the voice of God, a better alternative would be to let the two candidates settle the issue in a chariot race, or perhaps better, a dual. Even drawing lots would be fairer and more rational.

Because quantity is the sole criterion, any tactic whatsoever can be used to gain a vote. Issues are secondary, since very few people know what they are beyond the simple-minded propaganda statements that can fit on a placard. Even the self-anointed political pundits in the media understand this. They analyze events not in terms of whether they are true or false, just or unjust, competent or incompetent, effective or futile, intelligent or stupid, but only in terms of which candidate they favor in respect to the accumulation of votes. Ideas replaced by quantity.

Conclusion to follow.


\flrightit{Posted on 2012-11-15 by Cologero}

